\begin{theorem}[Toeplitz 证书塔的度量谱正规形与负谱平方律]\label{thm:xi-toeplitz-metric-spectrum-normal-form}
在有限缺陷口径下，取 \(N\) 足够大使 Toeplitz 证书矩阵 \(\mathbf{T}_N(C)\) 的负惯性稳定为 \(\kappa\)，并取 \(\mathsf H_{\kappa-1}\) 满秩，使得推论 \ref{cor:xi-toeplitz-scan-hankel-congruence-decomposition} 的合同分解成立：存在可逆矩阵 \(R_N\) 与某个 \(P_N\succeq 0\) 使
\begin{equation}\label{eq:xi-toeplitz-congruence-block-split}
R_N^{\ast}\,\mathbf{T}_N(C)\,R_N
=
\begin{pmatrix}
P_N & 0\\[2pt]
0 & -\,\mathsf H_{\kappa-1}^{\ast}\mathsf H_{\kappa-1}
\end{pmatrix}.
\end{equation}
定义正定审计度量
\begin{equation}\label{eq:xi-toeplitz-audit-metric-GN}
G_N:=R_N^{-*}R_N^{-1}\ \succ\ 0.
\end{equation}
则广义特征值问题
\begin{equation}\label{eq:xi-toeplitz-generalized-eig-GN}
\mathbf{T}_N(C)x=\lambda\,G_N x
\end{equation}
的广义谱等于块对角矩阵 \(\mathrm{diag}(P_N,-\mathsf H_{\kappa-1}^{\ast}\mathsf H_{\kappa-1})\) 的普通谱。
特别地，\eqref{eq:xi-toeplitz-generalized-eig-GN} 的全部负广义特征值恰为
\begin{equation}\label{eq:xi-negative-generalized-eigs-are-minus-sigmasq}
\lambda_i^-=-\sigma_i(\mathsf H_{\kappa-1})^2,\qquad i=1,\dots,\kappa,
\end{equation}
其中 \(\sigma_i(\mathsf H_{\kappa-1})\) 为 \(\mathsf H_{\kappa-1}\) 的奇异值（按任意顺序列举）。
并且相应的负广义特征子空间可取为
\begin{equation}\label{eq:xi-negative-generalized-eigenspace}
\mathcal N_N^-:=R_N\bigl(\{0\}\oplus\CC^{\kappa}\bigr)\subseteq \CC^{N+1}.
\end{equation}
对任意 \(y\in\CC^{\kappa}\) 令 \(x:=R_N(0\oplus y)\)，则严格恒等式成立
\begin{equation}\label{eq:xi-rayleigh-quotient-identities}
x^{\ast}\mathbf{T}_N(C)x=-\|\mathsf H_{\kappa-1}y\|_2^2,
\qquad
x^{\ast}G_N x=\|y\|_2^2.
\end{equation}
因此在 \(\mathcal N_N^-\setminus\{0\}\) 上的 Rayleigh 商由 \(\mathsf H_{\kappa-1}\) 的奇异值完全刻画：
\begin{equation}\label{eq:xi-rayleigh-negative-extrema-sigma}
\inf_{x\in\mathcal N_N^-\setminus\{0\}}
\frac{x^{\ast}\mathbf{T}_N(C)x}{x^{\ast}G_N x}
=-\sigma_{\max}(\mathsf H_{\kappa-1})^2,
\qquad
\sup_{x\in\mathcal N_N^-\setminus\{0\}}
\frac{x^{\ast}\mathbf{T}_N(C)x}{x^{\ast}G_N x}
=-\sigma_{\min}(\mathsf H_{\kappa-1})^2.
\end{equation}
\end{theorem}
\begin{proof}
由 \eqref{eq:xi-toeplitz-congruence-block-split} 令 \(x=R_N z\)，则
\[
R_N^{\ast}\mathbf{T}_N(C)R_N z
=
\begin{pmatrix}
P_N & 0\\
0 & -\,\mathsf H_{\kappa-1}^{\ast}\mathsf H_{\kappa-1}
\end{pmatrix}z.
\]
另一方面，由 \eqref{eq:xi-toeplitz-audit-metric-GN} 得
\(
G_N x=R_N^{-*}z
\)，从而
\(
R_N^{\ast}G_N x=z
\)。
因此 \(\mathbf{T}_N(C)x=\lambda G_N x\) 等价于
\[
\begin{pmatrix}
P_N & 0\\
0 & -\,\mathsf H_{\kappa-1}^{\ast}\mathsf H_{\kappa-1}
\end{pmatrix}z
=\lambda z,
\]
故广义谱与块对角普通谱一致。
负谱部分来自 \(-\mathsf H_{\kappa-1}^{\ast}\mathsf H_{\kappa-1}\)，其特征值为 \(-\sigma_i(\mathsf H_{\kappa-1})^2\)，得到 \eqref{eq:xi-negative-generalized-eigs-are-minus-sigmasq}。
取 \(z=(0\oplus y)\) 得到 \eqref{eq:xi-negative-generalized-eigenspace} 与 \eqref{eq:xi-rayleigh-quotient-identities}。
最后，由 Rayleigh 极值刻画立得 \eqref{eq:xi-rayleigh-negative-extrema-sigma}。证毕。
\end{proof}

\begin{corollary}[负谱乘积的 Hankel 行列式模平方不变量]\label{cor:xi-negative-spectrum-product-dethankel-square}
在定理 \ref{thm:xi-toeplitz-metric-spectrum-normal-form} 的设定下，\(\kappa\) 个负广义特征值满足
\begin{equation}\label{eq:xi-product-negative-eigs-dethankel}
\prod_{i=1}^{\kappa}|\lambda_i^-|
=
\det\!\bigl(\mathsf H_{\kappa-1}^{\ast}\mathsf H_{\kappa-1}\bigr)
=
\bigl|\det(\mathsf H_{\kappa-1})\bigr|^2.
\end{equation}
\end{corollary}
\begin{proof}
由 \eqref{eq:xi-negative-generalized-eigs-are-minus-sigmasq}，
\(
\prod_i|\lambda_i^-|=\prod_i\sigma_i(\mathsf H_{\kappa-1})^2
=\det(\mathsf H_{\kappa-1}^{\ast}\mathsf H_{\kappa-1})
=|\det(\mathsf H_{\kappa-1})|^2
\)。
证毕。
\end{proof}

\begin{theorem}[有限 Blaschke 扫描谱指纹下的 Hankel--Vandermonde 平方闭式]\label{thm:xi-dethankel-vandermonde-square-finite-blaschke}
设单侧扫描谱指纹满足指数分解
\begin{equation}\label{eq:xi-scan-fingerprint-exponential-sum}
u_n=4\pi\sum_{j=1}^{\kappa} w_j z_j^{n},\qquad n\ge 0,
\end{equation}
其中 \((w_1,\dots,w_\kappa)\in(\RR_{>0})^\kappa\)，且 \((z_1,\dots,z_\kappa)\in\mathbb{D}^{\kappa}\) 两两不同。
定义 Hankel 块
\(
\mathsf H_{\kappa-1}:=(u_{p+q})_{0\le p,q\le \kappa-1}
\)。
令 Vandermonde 矩阵 \(V\in\CC^{\kappa\times\kappa}\) 为
\[
V_{p+1,j}:=z_j^{p}\qquad(0\le p\le \kappa-1).
\]
则存在精确分解
\begin{equation}\label{eq:xi-hankel-vandermonde-factorization}
\mathsf H_{\kappa-1}=(4\pi)\,V\,\mathrm{diag}(w_1,\dots,w_\kappa)\,V^{\mathsf T}.
\end{equation}
从而
\begin{equation}\label{eq:xi-dethankel-vandermonde-square}
\det(\mathsf H_{\kappa-1})
=(4\pi)^{\kappa}\Bigl(\prod_{j=1}^{\kappa}w_j\Bigr)\Bigl(\prod_{1\le i<j\le\kappa}(z_j-z_i)^2\Bigr).
\end{equation}
结合推论 \ref{cor:xi-negative-spectrum-product-dethankel-square}，负广义谱乘积满足完全显式闭式
\begin{equation}\label{eq:xi-product-negative-eigs-explicit-vandermonde}
\prod_{i=1}^{\kappa}|\lambda_i^-|
=(4\pi)^{2\kappa}
\Bigl(\prod_{j=1}^{\kappa}w_j^{2}\Bigr)
\Bigl(\prod_{1\le i<j\le\kappa}|z_j-z_i|^{4}\Bigr).
\end{equation}
\end{theorem}
\begin{proof}
由 \eqref{eq:xi-scan-fingerprint-exponential-sum}，对任意 \(0\le p,q\le \kappa-1\) 有
\[
(\mathsf H_{\kappa-1})_{p+1,q+1}
=u_{p+q}
=4\pi\sum_{j=1}^{\kappa}w_j z_j^{p+q}
=4\pi\sum_{j=1}^{\kappa}V_{p+1,j}\,w_j\,V_{q+1,j},
\]
即得矩阵分解 \eqref{eq:xi-hankel-vandermonde-factorization}。
取行列式得到
\(
\det(\mathsf H_{\kappa-1})=(4\pi)^\kappa\det(V)^2\prod_j w_j
\)；
而 \(\det(V)=\prod_{i<j}(z_j-z_i)\) 为 Vandermonde 行列式，故得 \eqref{eq:xi-dethankel-vandermonde-square}。
最后将 \eqref{eq:xi-dethankel-vandermonde-square} 代入 \eqref{eq:xi-product-negative-eigs-dethankel} 得 \eqref{eq:xi-product-negative-eigs-explicit-vandermonde}。证毕。
\end{proof}

\begin{theorem}[Van Vleck 留数线性矩约束强制的残差能量下界]\label{thm:xi-vanvleck-residue-l2-lower-bound}
沿用命题 \ref{prop:xi-basepoint-scan-van-vleck-residue-moment01} 的设定。
令 \(p_k=s_k+\mathrm{i}t_k\in\CC\)（\(s_k,t_k\in\RR\)）与 \(r_k\in\CC\) 满足线性矩约束
\begin{equation}\label{eq:xi-vanvleck-residue-linear-moments}
\sum_{k=1}^{\kappa}\Re r_k=0,
\qquad
\sum_{k=1}^{\kappa}\Re(p_k r_k)= -\frac{\kappa(\kappa+1)}{2}.
\end{equation}
记
\[
S:=\sum_{k=1}^{\kappa}s_k,\qquad
Q:=\sum_{k=1}^{\kappa}|p_k|^2=\sum_{k=1}^{\kappa}(s_k^2+t_k^2).
\]
若 \(\kappa Q-S^2>0\)，则在约束 \eqref{eq:xi-vanvleck-residue-linear-moments} 下成立严格下界
\begin{equation}\label{eq:xi-vanvleck-residue-energy-lb}
\sum_{k=1}^{\kappa}|r_k|^2
\ \ge\
\frac{\kappa^3(\kappa+1)^2}{4(\kappa Q-S^2)}
\ =\
\frac{\kappa(\kappa+1)^2}{4\Bigl(\mathrm{Var}(s)+\frac1\kappa\sum_{k=1}^{\kappa}t_k^2\Bigr)},
\end{equation}
其中 \(\mathrm{Var}(s):=\frac1\kappa\sum_{k=1}^{\kappa}(s_k-\frac{S}{\kappa})^2\)。
并且达到等号的唯一最小能量解为
\begin{equation}\label{eq:xi-vanvleck-residue-energy-minimizer}
r_k
=
\frac{\kappa(\kappa+1)}{2(\kappa Q-S^2)}(S-\kappa s_k)
\;+\;
\mathrm{i}\,\frac{\kappa^2(\kappa+1)}{2(\kappa Q-S^2)}\,t_k.
\end{equation}
\end{theorem}
\begin{proof}
写 \(r_k=\alpha_k+\mathrm{i}\beta_k\)（\(\alpha_k,\beta_k\in\RR\)），则 \eqref{eq:xi-vanvleck-residue-linear-moments} 等价于
\[
\sum_{k=1}^{\kappa}\alpha_k=0,\qquad
\sum_{k=1}^{\kappa}(s_k\alpha_k-t_k\beta_k)= -\frac{\kappa(\kappa+1)}{2}.
\]
在这些线性约束下极小化二次型
\(
E(\alpha,\beta):=\sum_{k=1}^{\kappa}(\alpha_k^2+\beta_k^2)
=\sum_k|r_k|^2
\)。
取拉格朗日函数
\[
\mathcal L
=E(\alpha,\beta)
+\lambda_0\sum_{k=1}^{\kappa}\alpha_k
+\lambda_1\left(\sum_{k=1}^{\kappa}(s_k\alpha_k-t_k\beta_k)+\frac{\kappa(\kappa+1)}{2}\right).
\]
对 \((\alpha_k,\beta_k)\) 求偏导并令零得
\[
2\alpha_k+\lambda_0+\lambda_1 s_k=0,\qquad
2\beta_k-\lambda_1 t_k=0,
\]
即
\begin{equation}\label{eq:xi-vanvleck-residue-stationary}
\alpha_k=-\frac{\lambda_0+\lambda_1 s_k}{2},\qquad \beta_k=\frac{\lambda_1 t_k}{2}.
\end{equation}
由 \(\sum_k\alpha_k=0\) 得 \(\lambda_0=-(S/\kappa)\lambda_1\)。
将其代回第二条约束得到
\[
-\frac{\lambda_1}{2}\left(Q-\frac{S^2}{\kappa}\right)=-\frac{\kappa(\kappa+1)}{2},
\qquad\Longrightarrow\qquad
\lambda_1=\frac{\kappa^2(\kappa+1)}{\kappa Q-S^2}.
\]
再代回 \eqref{eq:xi-vanvleck-residue-stationary} 即得 \eqref{eq:xi-vanvleck-residue-energy-minimizer}。
最后，注意
\[
\sum_{k=1}^{\kappa}\Bigl(s_k-\frac{S}{\kappa}\Bigr)^2+\sum_{k=1}^{\kappa}t_k^2
=Q-\frac{S^2}{\kappa}
=\frac{\kappa Q-S^2}{\kappa},
\]
故
\[
\min E
=\frac{\lambda_1^2}{4}\left(Q-\frac{S^2}{\kappa}\right)
=\frac{\kappa^3(\kappa+1)^2}{4(\kappa Q-S^2)},
\]
得到 \eqref{eq:xi-vanvleck-residue-energy-lb}，并且极小解由严格凸性唯一。证毕。
\end{proof}

\begin{corollary}[极点几何有界性导出的统一能量下界]\label{cor:xi-vanvleck-residue-energy-uniform-lb-under-radius}
在定理 \ref{thm:xi-vanvleck-residue-l2-lower-bound} 的设定下，若另有半径约束
\[
|p_k|\le R\qquad(1\le k\le \kappa),
\]
则必有
\begin{equation}\label{eq:xi-vanvleck-residue-energy-uniform-lb}
\sum_{k=1}^{\kappa}|r_k|^2\ \ge\ \frac{\kappa(\kappa+1)^2}{8R^2}.
\end{equation}
\end{corollary}
\begin{proof}
由 \(|p_k|\le R\) 得 \(Q\le \kappa R^2\)，从而 \(\kappa Q-S^2\le \kappa Q\le \kappa^2R^2\)。
另一方面，
\(
\kappa Q-S^2=\kappa^2\mathrm{Var}(s)+\kappa\sum_k t_k^2\le 2\kappa Q
\)，故 \(\kappa Q-S^2\le 2\kappa^2R^2\)。
将其代入 \eqref{eq:xi-vanvleck-residue-energy-lb} 得到 \eqref{eq:xi-vanvleck-residue-energy-uniform-lb}。证毕。
\end{proof}

\begin{conjecture}[Poisson--KL 耗散上界与 Van Vleck 能量下界的证书矛盾闭环]\label{conj:xi-poisson-kl-vs-vanvleck-residue-energy-closed-loop}
在有限缺陷口径下，设 Toeplitz--扫描 Hankel 同构分解（推论 \ref{cor:xi-toeplitz-scan-hankel-congruence-decomposition}）与 Van Vleck 留数矩约束（命题 \ref{prop:xi-basepoint-scan-van-vleck-residue-moment01}）同时成立，并假设存在一族完全内部的可审计上界泛函 \(\mathcal{U}_m\ge 0\)（由 Poisson--KL 耗散链与 strictification 预算给出，见结论 \ref{con:xi-endpoint-defect-poisson-kl-second-order-energy-dominant-alignment} 的对齐口径），满足：
\begin{enumerate}
  \item 对每个分辨率层级 \(m\)，留数能量被上界控制：
  \[
  \sum_{k=1}^{\kappa}|r_k|^2\ \le\ \mathcal{U}_m;
  \]
  \item 当 \(m\to\infty\) 时有 \(\mathcal{U}_m\to 0\)。
\end{enumerate}
则必有 \(\kappa=0\)，从而点状缺陷消失，并由定理 \ref{thm:xi-disk-zero-pole-index-reduction} 得到 RH 的圆盘等价式成立。
\end{conjecture}

\endinput
